The metric below is the derived metric of the gravitational field $\hb^{-1}(e_{\tau}) = e_{0}$, $\hb^{-1}(e_{x^{i}}) = \sum_{l}\tau^{p_{l}}\theta^{l}(e_{x^{i}})\,e_{l}$ on flat spacetime (constant orthonormal physical frame $e_{0}, \dots, e_{3}$), and its Levi-Civita connection is the field's rotation-gauge field (Fact \claimref{synthesis.imported-gtg-dictionary}(ii)); the field equations $\mathcal{R}(a) = \Lambda a$ are $\mathrm{Ric} = \Lambda g$ for the derived metric. Each slice of constant $\tau$ carries the three-dimensional gauge field $\hb_{\tau}^{-1}(e_{x^{i}}) = \sum_{l}\tau^{p_{l}}\theta^{l}(e_{x^{i}})\,e_{l}$, whose physical vectors are the spacetime vectors orthogonal to $e_{0}$, the algebra $\Cl(0,3)$; its rotation-gauge field is fixed by the torsion equation, and $R^{(3)}$ is the scalar $\partial_{a}\cdot\mathcal{R}_{\tau}(a)$ of its store-sign curvature operator, which in this negative-definite signature is positive on the round sphere (dictionary C7: the store's sign and the round-sphere-positive one agree on the scalar).

(i) Fix $\tau$. The orthonormal frame of $\sum_{l}\tau^{2p_{l}}(\theta^{l})^{2}$ is $\hat e_{l} = \tau^{-p_{l}}E_{l}$, and since the weights are constant, $[\hat e_{m}, \hat e_{n}] = \hat c_{mn}{}^{l}\hat e_{l}$ with $\hat c_{mn}{}^{l} = \tau^{p_{l} - p_{m} - p_{n}}c_{mn}{}^{l}$. Decompose, as usual for three-dimensional frames, $\hat c_{ij}{}^{k} = \varepsilon_{ijl}\hat n^{lk} + \hat a_{i}\delta_{j}{}^{k} - \hat a_{j}\delta_{i}{}^{k}$ with $\hat n$ symmetric; then $\hat a_{i} = \tfrac12\hat c_{ij}{}^{j}$ and $\hat n^{lk} = \tfrac12\varepsilon^{lij}\hat c_{ij}{}^{k} - \varepsilon^{lik}\hat a_{i}$. The scalar curvature in this frame is
\[
R^{(3)} = 4\,\hat e_{i}(\hat a_{i}) - 6\,\hat a_{i}\hat a_{i} - \tfrac12\bigl(2\,\hat n^{kl}\hat n^{kl} - (\hat n^{kk})^{2}\bigr)
\]
(summed), the orthonormal-frame form of Ellis and MacCallum, which follows from the torsion equation. Let $e_{l}$ be the constant physical frame of the slice ($Q(e_{l}) = -1$), $\hat e_{l} = \hb_{\tau}(e_{l})$ its background vector fields and $\omega_{m} := \Omega(\hat e_{m})$. Since $\hb_{\tau}^{-1}(\hat e_{m}) = e_{m}$ is constant, the torsion equation for the fields $\hat e_{m}$, $D_{\hat e_{m}}e_{n} - D_{\hat e_{n}}e_{m} = \hb_{\tau}^{-1}([\hat e_{m}, \hat e_{n}])$, reads $\omega_{m}\times e_{n} - \omega_{n}\times e_{m} = \hat c_{mn}{}^{l}e_{l}$: nine linear equations for the nine components of $\omega_{1}, \omega_{2}, \omega_{3}$ with a constant invertible matrix, so each $\omega_{m}$ is linear in the $\hat c$ with constant coefficients. Expanding $\Omega(v) = v^{\mu}\Omega(e_{\mu})$ in the field strength gives $R(u \wedge v) = u\cdot\nabla\,\Omega(v) - v\cdot\nabla\,\Omega(u) + \Omega(u)\times\Omega(v) - \Omega([u, v])$ for vector fields $u, v$, so on the frame planes, in the store's sign,
\[
\mathcal{R}_{\tau}(e_{m}\wedge e_{n}) = -\bigl(\hat e_{m}(\omega_{n}) - \hat e_{n}(\omega_{m}) + \omega_{m}\times\omega_{n} - \hat c_{mn}{}^{l}\,\omega_{l}\bigr),
\]
and the scalar $\partial_{a}\cdot\mathcal{R}_{\tau}(a)$ is a polynomial in the nine $\hat c_{mn}{}^{l}$ (of degree two) and their twenty-seven frame derivatives $\hat e_{k}(\hat c_{mn}{}^{l})$ (of degree one). The check computes this polynomial exactly, with those thirty-six quantities as independent variables, and finds it equal, coefficient by coefficient, to the formula above written through the decomposition (whose $\hat n$ it confirms to be symmetric and to reproduce $\hat c$); only the six derivatives $\hat e_{i}(\hat c_{ij}{}^{j})$ occur, so no Jacobi identity is needed, and the left-invariant frame of the unit three-sphere ($\hat c_{ij}{}^{k} = \pm 2\varepsilon_{ijk}$) gives $6$ (check\_gtg\_ellis\_maccallum\_frame\_formula). The script also evaluates both on concrete fields: it solves the rotation-gauge field of the slice field from the torsion equation on exact second-order jets and finds the scalar $\partial_{a}\cdot\mathcal{R}_{\tau}(a)$ equal to the frame formula (and to the coordinate Ricci scalar of the jets of the metric), for three generic rational coframes and three weight triples each (check\_gtg\_slice\_and\_bianchi\_ii\_field\_strength). The statement itself is posed on the structure functions $c_{mn}{}^{l}$ of the coframe. The scalings are read off from $\sum_{l}p_{l} = 1$. The diagonal $\hat n^{kk} = \hat c_{ij}{}^{k}$, $(i, j, k)$ cyclic, scales as $\tau^{p_{k} - p_{i} - p_{j}} = \tau^{2p_{k} - 1}$. Every other ingredient carries one lower index equal to its upper index: $\hat c_{ij}{}^{j} = \tau^{-p_{i}}c_{ij}{}^{j}$, so $\hat a_{i}$ scales as $\tau^{-p_{i}}$; the off-diagonal $\hat n^{kl}$, $k \neq l$, is built from $\hat c_{ij}{}^{k}$ with $\{i, j\}$ the complement of $l$ (so $k \in \{i, j\}$) and from $\hat a$, all scaling as $\tau^{-p_{m}}$ with $m$ the index other than $k, l$; and $\hat e_{i}(\hat a_{i}) = \tau^{-2p_{i}}E_{i}(\tfrac12 c_{ij}{}^{j})$. Multiplying by $\tau^{2}$: $(\hat n^{kk})^{2}$ gives $\tau^{4p_{k}}$; the cross terms $\hat n^{kk}\hat n^{ll}$ of $(\mathrm{tr}\,\hat n)^{2}$ give $\tau^{2p_{k} + 2p_{l}} = \tau^{2 - 2p_{m}}$; every remaining term gives $\tau^{2 - 2p_{m}}$ for its index $m$. The coefficient of $\tau^{4p_{k}}$ is $-\tfrac12(2 - 1)(c_{ij}{}^{k})^{2} = -\tfrac12 h_{k}^{2}$, where $h_{k} = -c_{ij}{}^{k}$ because $\theta^{k}\wedge d\theta^{k}(E_{1}, E_{2}, E_{3}) = d\theta^{k}(E_{i}, E_{j}) = -\theta^{k}([E_{i}, E_{j}])$ for $(i, j, k)$ cyclic. The coefficients $Q_{m}$ collect the rest and contain no weights. The script computes $h_{k}$ as the trivector part of the geometric product of $\theta^{k}$ with $d\theta^{k}$ in $\Cl(3,0)$, pins the values, confirms $h_{k} = -c_{ij}{}^{k}$, and confirms the support identity as an identity in the weights: with $a_{l} = \tau^{p_{l}}$ written as independent rational $a_{l}$, $R^{(3)} + \tfrac12\sum_{k}h_{k}^{2}a_{k}^{2}/(a_{i}a_{j})^{2} = \sum_{m}Q_{m}/a_{m}^{2}$ with $Q$ fitted at three weight triples and confirmed at two more, for each coframe, while dropping the helicity term breaks it.

(ii) On the Kasner circle an exponent vanishes only at the permutations of $(1, 0, 0)$ and some exponent is nonpositive (\claimref{synthesis.unit-boost-residue-needs-nonpositive-exponent}(2); in the parametrisation $p(u) = (-u, 1 + u, u(1 + u))/(1 + u + u^{2})$, $u \geq 1$); away from those points exactly one exponent is negative and the others lie in $(0, 1)$. So $4p_{1} < 0$, while $4p_{2}, 4p_{3} > 0$ and $2 - 2p_{m} \geq 2 - 2p_{3} > 0$. If $\theta(\tau)$ converges in $C^{2}$, the structure functions and their first derivatives converge, so $h_{k}(\tau, x)$ and $Q_{m}(\tau, x)$ are bounded, and (i) gives $\tau^{2}R^{(3)} = -\tfrac12 h_{1}^{2}\tau^{4p_{1}} + O(\tau^{\delta})$ with $\delta = \min(4p_{2}, 2 - 2p_{3})$. The first term tends to zero if and only if $\tau^{2p_{1}}h_{1} \to 0$; for a coframe independent of $\tau$, $\tau^{4p_{1}} \to \infty$, so if and only if $h_{1}(x) = 0$. The script pins $\delta$ at $u = 1, \tfrac32, 2, \tfrac72$ (values $\tfrac23, \tfrac8{19}, \tfrac27, \tfrac8{67}$), confirms that $4p_{1}$ is the only nonpositive exponent there, and that at $(1, 0, 0)$ an exponent is zero, where the statement degenerates. Off the circle, on the part of the scalar disk ($\sum_{l}p_{l} = 1$, $\sum_{l}p_{l}^{2} < 1$) where every $p_{l}$ is in $(0, 1)$ (nonempty: it contains $(\tfrac13, \tfrac13, \tfrac13)$; the disk point $(-\tfrac16, \tfrac56, \tfrac13)$ shows an exponent can be negative elsewhere in it), all exponents $4p_{k}$ and $2 - 2p_{m}$ in (i) are positive, so $\tau^{2}R^{(3)} \to 0$ for every coframe converging in $C^{2}$ (filed alongside, synthesis.scalar-kasner-wall-is-reduced-five-dimensional-vacuum-wall).

(iii) Under the criterion of \claimref{synthesis.nondiagonal-kasner-residue-weighted-shear}(ii) with residue $P$, $\tau K = P + o(1)$, and with $\tau\partial_{\tau}(\tau K) \to 0$ the $00$ component of the field equations at order $\tau^{-2}$ is $\mathrm{tr}P = \mathrm{tr}P^{2}$, as in the proof of clause (iv) there, and it holds on the circle. The Hamiltonian constraint, the $\gamma_{0}$ component of the Einstein map, $-2\gamma_{0}\cdot\mathcal{G}(\gamma_{0}) = (\mathrm{tr}K)^{2} - \mathrm{tr}(K^{2}) + s\,R^{(3)} = 2\Lambda$ with $s = \pm 1$ (the Gauss identity of the slices as stated in the proof of \claimref{synthesis.nondiagonal-kasner-residue-weighted-shear}(iv), quoted there and not derived in geometric algebra for curved slices; the sign $s$ does not matter here, since only the vanishing of the limit of $\tau^{2}R^{(3)}$ enters), multiplied by $\tau^{2}$, reads at order one $(\mathrm{tr}P)^{2} - \mathrm{tr}P^{2} + \lim\tau^{2}R^{(3)} = 0$; on the circle $(\mathrm{tr}P)^{2} = \mathrm{tr}P^{2} = 1$, so it holds if and only if $\tau^{2}R^{(3)} \to 0$, which by (ii) is $\tau^{2p_{1}}h_{1} \to 0$; $\Lambda$ enters only at order $\tau^{2}$. If the limiting coframe has $h_{1}(x) \neq 0$, then $\tau^{2}R^{(3)} = -\tfrac12 h_{1}^{2}\tau^{4p_{1}}(1 + o(1)) \to -\infty$ at $x$, so $\tau^{2}$ times the Hamiltonian constraint diverges like $\tau^{4p_{1}}$ (its curvature term is of order $\tau^{4p_{1} - 2}$, more singular than $\tau^{-2}$), and no solution of the field equations has this form there. The momentum constraints and the spatial components are not used. The script checks that the residue part $(\mathrm{tr}P)^{2} - \mathrm{tr}P^{2}$ vanishes at the sample points of the circle and not at an off-circle control with $\sum p = 1$.

(iv) Take coordinates $(T, x, y, z)$ with $\theta^{1} = dx - z\,dy$, $\theta^{2} = dy$, $\theta^{3} = dz$, so $d\theta^{1} = dy\wedge dz = \theta^{2}\wedge\theta^{3}$. The metric is invariant under the left translations of the Heisenberg group, which act transitively on each slice, so it suffices to compute at $z = 0$. With $\alpha_{l}$, $\alpha_{l}'$, $\alpha_{l}''$ as independent rationals (the jets of $a_{l}$ follow), the script builds the gauge field $\hb^{-1}(e_{T}) = N e_{0}$, $\hb^{-1}(e_{x}) = a_{1}e_{1}$, $\hb^{-1}(e_{y}) = a_{2}e_{2} - z\,a_{1}e_{1}$, $\hb^{-1}(e_{z}) = a_{3}e_{3}$ on exact jets in $(T, z)$ and solves its rotation-gauge field from the torsion equation; at every $z$ it is
\[
\Omega(e_{T}) = 0, \qquad \Omega(e_{x}) = -\frac{\alpha_{1}'}{a_{2}a_{3}}\,e_{0}\wedge e_{1} + \frac{a_{1}^{2}}{2a_{2}a_{3}}\,e_{2}\wedge e_{3}, \qquad \Omega(e_{z}) = -\frac{\alpha_{3}'}{a_{1}a_{2}}\,e_{0}\wedge e_{3} - \frac{a_{1}}{2a_{2}}\,e_{1}\wedge e_{2},
\]
\[
\Omega(e_{y}) = z\,\frac{\alpha_{1}'}{a_{2}a_{3}}\,e_{0}\wedge e_{1} - \frac{\alpha_{2}'}{a_{1}a_{3}}\,e_{0}\wedge e_{2} + \frac{a_{1}}{2a_{3}}\,e_{1}\wedge e_{3} - z\,\frac{a_{1}^{2}}{2a_{2}a_{3}}\,e_{2}\wedge e_{3}
\]
(primes are $d/dT$; the check compares the solved field with this closed form, value and first derivatives in $T$ and $z$). The coordinate Ricci components $R_{\mu\nu} = -g_{\mu}\cdot\mathcal{R}(g_{\nu})$ (round-sphere-positive sign, G2) follow from it by direct differentiation, identically in the $\alpha_{l}(T)$; the script reads them off the Ricci map and finds, exactly at two generic jets (and the same values from the Christoffel symbols of the metric's jets), $R_{TT} = -\tfrac12(E_{1} + E_{2} + E_{3}) + \tfrac12 C$, $N^{2}R(e_{l}, e_{l}) = \tfrac12E_{l}$ and vanishing off-diagonal components, with $E_{1} = 2\alpha_{1}'' + a_{1}^{4}$, $E_{2,3} = 2\alpha_{2,3}'' - a_{1}^{4}$ and $C = 4(\alpha_{1}'\alpha_{2}' + \alpha_{2}'\alpha_{3}' + \alpha_{3}'\alpha_{1}') - a_{1}^{4}$ (the type II case of the standard Bianchi IX equations, Landau and Lifshitz, cited from memory); it also checks that the Ricci map $\mathcal{R}(a)$ vanishes at two jets solving the four equations and not when one is perturbed. So vacuum holds if and only if $E_{l} = 0$ and $C = 0$. The first is Liouville's equation $2\alpha_{1}'' = -e^{4\alpha_{1}}$; multiplying by $\alpha_{1}'$ gives the first integral $\alpha_{1}'^{2} + \tfrac14 e^{4\alpha_{1}} = v^{2}$ with $v > 0$, whose solutions are exactly $\alpha_{1} = \tfrac12\log(2v) - \tfrac12\log\cosh(2v(T - T_{0}))$: with $s = \tanh(2v(T - T_{0}))$, $\alpha_{1}' = -vs$, $\alpha_{1}'' = -2v^{2}(1 - s^{2})$, $e^{4\alpha_{1}} = 4v^{2}(1 - s^{2})$, identities the script checks at rational $s$. Then $(\alpha_{1} + \alpha_{2})'' = (\alpha_{1} + \alpha_{3})'' = 0$, so $\alpha_{2}' = -\alpha_{1}' + B$, $\alpha_{3}' = -\alpha_{1}' + C_{0}$, and $C = 0$ reads $4(BC_{0} - \alpha_{1}'^{2}) = e^{4\alpha_{1}}$, that is $BC_{0} = v^{2}$ by the first integral. As $T \to \pm\infty$, $s \to \pm1$, $e^{4\alpha_{1}} \to 0$ and the velocities tend to $(-v, B + v, C_{0} + v)$ and $(v, B - v, C_{0} - v)$; with $BC_{0} = v^{2}$ each normalised triple lies on the Kasner circle ($\sum_{i<j}$ of the velocity products is $BC_{0} - v^{2} = 0$), and the field tends to the corresponding Kasner epochs with $T = \log t + \mathrm{const}$ in the proper time of each. Let $S_{\pm}$ be the sums; with $p = (-v, B + v, C_{0} + v)/S_{+}$, $B + v = p_{2}S_{+}$, $C_{0} + v = p_{3}S_{+}$ and $v = -p_{1}S_{+}$ give $S_{-} = S_{+}(1 + 2p_{1})$, and $1 + 2p_{1} \geq \tfrac13$ on the circle, so for $S_{+} > 0$ both sums are positive, the volume grows with $T$ at both ends and the singularity is at $T \to -\infty$. Then $p_{1} = -v/S_{+} < 0$ and
\[
p' = \frac{(v,\ B - v,\ C_{0} - v)}{S_{-}} = \frac{(-p_{1},\ p_{2} + 2p_{1},\ p_{3} + 2p_{1})}{1 + 2p_{1}},
\]
the BKL transition law; on $p(u)$ it is $p(u - 1)$ up to order for $u \geq 2$ and $p(1/(u - 1))$ for $1 < u < 2$, as recalled in \claimref{synthesis.bkl-bounces-no-framewise-invariant}. In $p'$ the first weight is $-p_{1}/(1 + 2p_{1}) > 0$ and $p_{2} + 2p_{1} = (1 - u)/(1 + u + u^{2}) \leq 0$, so the nonpositive weight is on axis $2$ or $3$, where $h_{2} = h_{3} = 0$, and by (ii) the curvature of that epoch's slices is subleading. The script checks all of this at three $(v, B)$ and five values of $s$, and the $u$-form at four values of $u$.
